\chapter*{符号说明}
\markboth{符号说明}{符号说明}
\phantomsection\addcontentsline{toc}{chapter}{符号说明}
\SourcePage{20}
下面列出本书将用到的部分符号。

\begingroup\linespread{1.18}\selectfont
\setlength{\LTpre}{4pt}\setlength{\LTpost}{4pt}\renewcommand{\arraystretch}{0.96}
\begin{longtable}{@{}p{0.38\textwidth}p{0.57\textwidth}@{}}
\toprule
\multicolumn{2}{@{}l}{\textbf{一般符号}}\\
\midrule
$M$ & 复流形或辛流形\\
$X$ 与 $V$ & 代数簇与 Calabi--Yau 簇\\
$\Theta_X$ 与 $T_X$ & $X$ 的切层与切丛\\
$\omega_X$ 与 $K_X$ & $X$ 的典范丛与典范除子\\
$\widehat\Omega_X^p$ & 轨形 $X$ 上的 Zariski $p$-形式\\
$M(X)$ & $X$ 上有效一维圈的 Mori 锥\\
$A_k(X)$ & $X$ 的第 $k$ 个 Chow 群\\
$\langle a,b\rangle$ 或 $g(a,b)$ & 上同调上的杯积配对\\
$\mathbb P^n$ 与 $\mathbb P(q_0,\ldots,q_n)$ & 射影空间与加权射影空间\\
\midrule
\multicolumn{2}{@{}l}{\textbf{环面簇与多胞形（\hyperref[chap:toric]{第 3 章}）}}\\
\midrule
$\Sigma$ 与 $\Sigma(1)$ & 扇及其一维锥\\
$\Delta$ 与 $\Delta^\circ$ & 多胞形及其极多胞形（或对偶多胞形）\\
$L(\Delta)$ & 指数属于 $\Delta$ 的 Laurent 多项式\\
$X_\Sigma$ 与 $\mathbb P_\Delta$ & 分别由扇 $\Sigma$ 和多胞形 $\Delta$ 确定的环面簇\\
$\operatorname{cpl}(\Sigma)$ & $\Sigma$ 上凸支撑函数构成的锥\\
\midrule
\multicolumn{2}{@{}l}{\textbf{Hodge 理论与 Yukawa 耦合（\hyperref[chap:hodge-yukawa]{第 5 章}）}}\\
\midrule
$\nabla=\nabla^{\mathrm{GM}}$ & Gauss--Manin 联络\\
$\mathcal F^\bullet$ 与 $\mathcal W_\bullet$ & Hodge 滤过与权滤过\\
$\mathcal H$ 与 $\mathcal H_{\mathbb Z}$ & Hodge 丛 $\mathcal F^0$ 及其中的整数子层\\
$T_j$ 与 $N_j$ & 单值性变换及其对数\\
$\left\langle\frac{\partial}{\partial t_i},\frac{\partial}{\partial t_j},\frac{\partial}{\partial t_k}\right\rangle=Y_{ijk}$ & 归一化 Yukawa 耦合，或 B 模型关联函数\\
\midrule
\multicolumn{2}{@{}l}{\SourcePage{21}\textbf{复模与 Kähler 模（\hyperref[chap:moduli]{第~\ref*{chap:moduli} 章}）}}\\
\midrule
$\mathcal M$ 与 $\overline{\mathcal M}$ & 复模空间及其紧化\\
$\mathcal M_{\mathrm{poly}}$ 与 $\overline{\mathcal M}_{\mathrm{poly}}$ & 多项式模空间及其紧化\\
$\mathcal M_{\mathrm{simp}}$ 与 $\overline{\mathcal M}_{\mathrm{simp}}$ & 简化模空间及其紧化\\
$K(V)$ 与 $\overline{K(V)}_{\mathbb C}$ & Kähler 锥与复化 Kähler 空间\\
$\mathcal{KM}$ 与 $\overline{\mathcal{KM}}$ & Kähler 模空间及其紧化\\
$\mathcal{KM}_{\mathrm{toric}}$ 与 $\overline{\mathcal{KM}}_{\mathrm{toric}}$ & 环面 Kähler 模空间及其紧化\\
\midrule
\multicolumn{2}{@{}l}{\textbf{Gromov--Witten 不变量与量子上同调（第 \hyperref[chap:gromov-witten]{\ref*{chap:gromov-witten}}、\hyperref[chap:quantum-cohomology]{\ref*{chap:quantum-cohomology}} 章）}}\\
\midrule
$\overline M_{g,n}(X,\beta)$ 与 $\overline{\mathcal M}_{g,n}(X,\beta)$ & 类为 $\beta$、亏格为 $g$、带 $n$ 个标记点的稳定映射的粗模空间与精细模叠\\
$[\overline{\mathcal M}_{g,n}(X,\beta)]^{\mathrm{virt}}$ & 虚拟基本类\\
$I_{g,n,\beta}(\alpha_1,\ldots,\alpha_n)$ 与 $\langle I_{g,n,\beta}\rangle(\alpha_1,\ldots,\alpha_n)$ & Gromov--Witten 类与 Gromov--Witten 不变量\\
$n_\beta$ 与 $n_d$ & 一般情形及五次三维簇的瞬子数\\
$N_\beta=\langle I_{0,0,\beta}\rangle$ 与 $N_d=\langle I_{0,0,d}\rangle$ & 一般情形及五次三维簇的零标记点、亏格零 Gromov--Witten 不变量\\
$*_{\mathrm{small}}$ 与 $*$ & 小量子乘积与大量子乘积（小量子乘积有时也记为 $*$）\\
$\langle a,b,c\rangle$ & 三点函数，或 A 模型关联函数\\
$\Phi(\gamma)$ & 亏格零 Gromov--Witten 势\\
$\nabla=\nabla^{\mathrm{GW}}$ & A 模型联络\\
\midrule
\multicolumn{2}{@{}l}{\textbf{Chern 类与等变上同调（\hyperref[chap:localization]{第~\ref*{chap:localization} 章}）}}\\
\midrule
$c_i(\mathcal E)$ 与 $c_i^T(\mathcal E)$ & $\mathcal E$ 的 Chern 类与等变 Chern 类\\
$\operatorname{Euler}(\mathcal E)$ 与 $\operatorname{Euler}_T(\mathcal E)$ & $\mathcal E$ 的 Euler 类与等变 Euler 类\\
$\int_{X_T}$ & 等变积分\\
$\lambda_0,\ldots,\lambda_n$ 与 $\hbar$ & $H^*(BT)$ 与 $H^*(B\mathbb C^*)$ 的环生成元；当 $G=\mathbb C^*\times T$、$T=(\mathbb C^*)^{n+1}$ 时，它们生成 $H^*(BG)$\\
$H$ 与 $p$ & 超平面类与等变超平面类\\
$q_0,\ldots,q_n$ & $\mathbb P^n$ 上标准 $T$ 作用的不动点\\
\midrule
\multicolumn{2}{@{}l}{\SourcePage{22}\textbf{引力关联数与 $J$ 函数（\hyperref[chap:quantum-differential-equations]{第~\ref*{chap:quantum-differential-equations} 章}）}}\\
\midrule
$\langle\tau_{d_1}\gamma_1,\ldots,\tau_{d_n}\gamma_n\rangle_{g,\beta}$ & 引力关联数\\
$\langle\gamma_1,\ldots,\gamma_n\rangle_{g,\beta}$ & $\langle I_{g,n,\beta}\rangle(\gamma_1,\ldots,\gamma_n)$ 的另一种记法\\
$\langle\!\langle\tau_{d_1}\gamma_1,\ldots,\tau_{d_n}\gamma_n\rangle\!\rangle_g$ & 亏格为 $g$ 的引力耦合\\
$\Phi_g^{\mathrm{grav}}(\gamma)$ & 亏格为 $g$ 的引力 Gromov--Witten 势\\
$*_{\mathrm{grav}}$ & 引力量子乘积\\
$\nabla^g$ 与 $\widetilde\nabla^g$ & Givental 联络及其对偶联络\\
$s_j$ & Givental 联络的平坦截面\\
$\left\langle\!\left\langle\dfrac{T_a}{\hbar-c},T_j\right\rangle\!\right\rangle_0$ & $\displaystyle\sum_{n=0}^{\infty}\hbar^{-(n+1)}\langle\!\langle\tau_n T_a,T_j\rangle\!\rangle_0$ 的符号记法\\
$J=J_X$ & 光滑簇 $X$ 的 Givental $J$ 函数\\
\midrule
\multicolumn{2}{@{}l}{\textbf{镜像定理（\hyperref[chap:mirror-theorem]{第~\ref*{chap:mirror-theorem} 章}）}}\\
\midrule
$\mathcal V=\bigoplus_{i=1}^{\ell}\mathcal O_{\mathbb P^n}(a_i)$ & 镜像定理中使用的向量丛\\
$\mathcal V_d$ 与 $\mathcal V_{d,k}$ & 由 $\mathcal V$ 诱导的 $\overline M_{0,0}(\mathbb P^n,d)$ 与 $\overline M_{0,k}(\mathbb P^n,d)$ 上的丛\\
$M_d$ 与 $N_d$ & $\overline M_{0,0}(\mathbb P^1\times\mathbb P^n,(1,d))$ 与 $\mathbb P(H^0(\mathbb P^1,\mathcal O_{\mathbb P^1}(d))^{n+1})$ 的简记\\
$p_{i,r}$ & $N_d$ 上 $G=\mathbb C^*\times T$ 作用的不动点\\
$\kappa$ & $H_G^*(N_d)$ 中的等变超平面类\\
$\mathcal R_T$ 与 $\mathcal R_G$ & $H^*(BT)$ 与 $H^*(BG)$ 的分式域\\
$\widehat P$ 与 $\widehat Q$ & 重要的 Euler 数据\\
$HG[B](t)$ & 集合 $B=\{B_d\}$ 所确定的上同调值函数\\
$\Psi(t)$ & 镜像映射\\
$I_{\mathcal V}$ 与 $J_{\mathcal V}$ & Givental 的上同调值函数\\
$\mathcal V'_{d,k,1}$ & 定义 $J_{\mathcal V}$ 所用的 $\mathcal V_{d,k}$ 的某个子丛\\
$*_X$ & 由完全交 $X\subset\mathbb P^n$ 决定的 $H^*(\mathbb P^n)$ 上的修正量子乘积\\
$\overline I_{\mathcal V}$ & 量子超平面截面原理中使用的上同调值函数\\
$I_T$ 与 $J_T$ & $I_{\mathcal V}$ 与 $J_{\mathcal V}$ 的等变版本\\
$Z_{i,\mathcal V}$ & 对 $0\le i\le n$ 定义的一组函数，它们唯一确定 $J_T$\\
\bottomrule
\end{longtable}
\endgroup

参考文献的引用约定，在\hyperlink{source:460}{参考文献表}的开头说明。
